\documentclass[12pt]{report}
\usepackage[brazilian]{babel}
\usepackage{csquotes}
\usepackage{makeidx}
\usepackage{glossaries}
\usepackage[
  perfil=relatorio,
  impressao=frente-e-verso,
  citacoes=autor-data
]{abntex3}

\makeindex
\makeglossaries
\newglossaryentry{reprodutibilidade}{
  name={reprodutibilidade},
  description={capacidade de repetir um procedimento documentado}
}
\ABNTEXsigla{API}{Interface de programação de aplicações}
\ABNTEXsimbolo{H}{Hipótese experimental}
\addbibresource{referencias-exemplo.bib}

\ABNTEXsetup{
  titulo={Avaliação experimental de formatos documentais},
  subtitulo={relatório científico completo para auditoria},
  autoria={Carla Souza e Daniel Alves},
  instituicao={Laboratório Exemplo de Ciência da Informação},
  volume={1},
  local={Recife},
  data={2026}
}

\ABNTEXrelatorioSetup{
  projeto-programa-plano={Projeto Formatos Abertos},
  numero-relatorio={RC-2026-001},
  codigo-identificacao={LECI-RC-2026-FA-001},
  classificacao-seguranca={acesso público},
  tipo-relatorio={relatório científico},
  resumo={Este relatório descreve um ensaio controlado de preservação,
    registra o método, sintetiza os resultados e apresenta recomendações
    para estudos posteriores.},
  palavras-chave={formatos abertos, reprodutibilidade, preservação},
  instituicao-endereco={Rua da Pesquisa, 10, Recife, PE},
  instituicao-patrocinadora={Fundação Exemplo de Pesquisa},
  patrocinadora-endereco={Avenida da Ciência, 50, Brasília, DF},
  issn={0000-0000},
  edicao={1},
  numero-paginas={a preencher após a composição},
  numero-classificacao={CDU 000.00},
  tiragem={digital},
  preco={distribuição gratuita},
  distribuidor={Laboratório Exemplo de Ciência da Informação},
  observacoes={Dados e nomes exclusivamente demonstrativos},
  possui-citacoes=sim
}

\ABNTEXrelatoriomembroequipe
  {Carla Souza}{Coordenação científica}{Laboratório Exemplo}
\ABNTEXrelatoriomembroequipe
  {Daniel Alves}{Execução experimental}{Laboratório Exemplo}

\begin{document}
\ABNTEXrelatorioexterna
\ABNTEXrelatoriocapa
\ABNTEXlombada[RC-2026-001]

\ABNTEXrelatoriopretextual
\ABNTEXrelatoriofolhaderosto
\ABNTEXrelatorioequipe
\ABNTEXerrata
  {SOUZA, Carla; ALVES, Daniel. Avaliação experimental de formatos
   documentais. Recife: Laboratório Exemplo, 2026.}
  {Não há correções nesta versão demonstrativa.}
\ABNTEXagradecimentos{À equipe de apoio do laboratório.}
\ABNTEXrelatorioresumo
\ABNTEXlistadeilustracoes
\ABNTEXlistadetabelas
\ABNTEXlistadesiglas
\ABNTEXlistadesimbolos
\ABNTEXsumario

\ABNTEXrelatoriotextual
\ABNTEXrelatoriointroducao{%
  O estudo avalia formatos documentais e registra os objetivos e as razões
  para a elaboração do relatório.
}
\ABNTEXrelatoriodesenvolvimento{%
  O ensaio compara propriedades observáveis e mantém registro
  \gls{reprodutibilidade}\index{reprodutibilidade}
  das decisões. A normalização bibliográfica permanece delegada ao
  \texttt{biblatex-abnt} \parencite{associacao2025}.

  \begin{figure}[htbp]
    \centering
    \rule{0.55\textwidth}{3cm}
    \caption{Fluxo experimental demonstrativo}
    \ABNTEXfonte{Elaborada pela equipe (2026).}
  \end{figure}

  \begin{table}[htbp]
    \centering
    \caption{Resultados sintéticos}
    \begin{tabular}{lc}
      \hline
      Critério & Resultado \\
      \hline
      Integridade & Verificada \\
      Reprodutibilidade & Verificada \\
      \hline
    \end{tabular}
    \ABNTEXfonte{Dados demonstrativos (2026).}
  \end{table}
}
\ABNTEXrelatorioconsideracoesfinais{%
  Os procedimentos são reproduzíveis no conjunto observado, mas resultados
  externos dependem de novos ensaios e revisão por especialistas.
}

\ABNTEXrelatoriopostextual
\ABNTEXbibliografia
\ABNTEXglossario
\ABNTEXappendix{Protocolo elaborado pela equipe}
Sequência demonstrativa de procedimentos do experimento.
\ABNTEXannex{Especificação externa}
Documento externo usado como referência no ensaio.
\ABNTEXindice[Índice de assuntos]
\ABNTEXrelatorioformulario
\end{document}
